\documentclass[10pt,a4paper]{amsart}
\usepackage[margin=19mm]{geometry}
\usepackage{amsmath,amssymb,mathtools}
\usepackage[scr=rsfs,cal=euler]{mathalfa}
\usepackage{xfrac}
\usepackage[unicode,hidelinks,bookmarks=false]{hyperref}
\newcommand{\ie}{i.e.\ }
\newcommand{\cf}{cf.\ }
\newcommand{\resp}{respectively\ }
\newcommand{\Subsection}{\subsection}
\providecommand{\nobreakdash}{\nobreak\hbox{-}}
\setlength{\emergencystretch}{2em}
\multlinegap0pt
\usepackage{bm}
\usepackage{bbm}
\usepackage{dsfont}
\usepackage{xspace}
\usepackage{cancel}
\usepackage{mathtools}
\usepackage{amsthm}
\makeatletter
\@ifundefined{theorem}{\newtheorem{theorem}{Theorem}}{}
\@ifundefined{proposition}{\newtheorem{proposition}[theorem]{Proposition}}{}
\makeatother
\DeclareMathOperator{\GL}{GL}
\DeclareMathOperator{\SO}{SO}
\newcommand{\II}{\mathbf{I}_{d \times d}}
\newcommand{\R}{\mathbb{R}}
\newcommand{\mc}[1]{\mathcal{#1}}
\providecommand\Lpl; %
\renewcommand{\Lpl}{\Delta}
\providecommand\X; %
\renewcommand{\X}{X}
\providecommand\a; %
\renewcommand{\a}{\alpha}
\providecommand\x; %
\renewcommand{\x}{x}
\title[Yang-Mills Meets Data]{Yang-Mills Meets Data}
\author{Jonas Cassel, Fabio Schlindwein, Peter Albers, Christoph Schn\"orr}
\date{Source: arXiv:2510.19431v1 (CC BY 4.0)}
\begin{document}
\maketitle

\noindent\emph{Source and licence.} Yang-Mills Meets Data. Version arXiv:2510.19431v1 (CC BY 4.0), distributed under the Creative Commons Attribution 4.0 International licence, as recorded in the arXiv licence metadata for this exact version and shown on the abstract page https://arxiv.org/abs/2510.19431v1. Full rights notice: licenses/arxiv-2510.19431-CC-BY-4.0.txt.\par\smallskip

\noindent\textit{Editorial scope.} Counted unit: the Proposition whose source label is \texttt{prop:loc-lpl} in the article (source0.tex lines 866--880, source label \texttt{prop:loc-lpl}), delivered together with the article's own complete original proof of it (source0.tex lines 882--985, one \texttt{proof} environment of 104 lines). No mathematics is written, repaired or re-derived: statement and proof are the article's own text line for line. Required context retained uncounted: eq:consis (lines 365--367); eq:def-lalpha (lines 659--661); eq:def-yi (lines 846--848); eq:def-y-diag (lines 850--852); eq:ass-g (lines 856--858). Every definition, display and label of those lines is retained unchanged. Omitted from this package and not redistributed: 1--364, 368--658, 662--845, 849--849, 853--855, 859--865, 881--881, 986--2250. Nothing inside a retained excerpt is omitted or altered: every retained body is delivered exactly as the article's lines are written, so no omission receipt of any kind accompanies this package. Citations: the retained excerpts carry no citation command at all, so no bibliography entry of the article is transcribed and no reference is redistributed. Reference adaptation: none was needed and none was made; every reference inside the delivered unit resolves inside it, so no unresolved reference or citation is produced. Printed slips kept literal: no printed word, symbol, hypothesis or number of the counted statement or of its proof was altered. Layout and class: the article's own class is \texttt{amsart} with packages including libertine, amsmath, amssymb, amsthm, mathtools, tikz-cd, float, wrapfig, caption, subcaption, aligned-overset, inputenc, hyperref, wasysym; the wrapper is the lane's pinned \texttt{amsart} editorial wrapper at 10pt on A4, which restates the theorem-like environments the retained text needs and adds the article's own macro definitions that the retained text needs (\texttt{\textbackslash GL}, \texttt{\textbackslash II}, \texttt{\textbackslash Lpl}, \texttt{\textbackslash R}, \texttt{\textbackslash SO}, \texttt{\textbackslash X}, \texttt{\textbackslash a}, \texttt{\textbackslash mc}, \texttt{\textbackslash x}), with extra wrapper packages (bm, bbm, dsfont, xspace, cancel, mathtools). The delivered numbering is the unit's own. Assets: no retained span contains a figure environment, a graphics-inclusion command, a PDF-page inclusion, a table or a TikZ picture; no image file is copied into this package, the package declares no asset dependency and no image file is redistributed with it. Licence: version arXiv:2510.19431v1 of this article is distributed under the Creative Commons Attribution 4.0 International licence, as recorded in the arXiv licence metadata for this exact version and shown on the abstract page https://arxiv.org/abs/2510.19431v1. A full rights notice is delivered at licenses/arxiv-2510.19431-CC-BY-4.0.txt. Characters: every retained excerpt is pure ASCII, so the delivered engine, XeLaTeX with its default Latin Modern text font, reports no missing character.
		\begin{equation}\label{eq:consis}
			F(gh^{-1},hx) = h F(g,\x ), \quad \forall h \in \GL(\X).
		\end{equation}

			\begin{equation}\label{eq:def-lalpha}
				(\Lpl_{\eta,\kappa,\alpha}x)_i = -\sum_{j \in \nu} \frac{\kappa_{ij}}{Z_i}(\a_{ij} \x_j - \x_i) \in \R^d, \qquad i \in \nu.
			\end{equation}

	\begin{equation}\label{eq:def-yi}
		\mc{F}([\xi^{i},x]) = [\xi^{i},y^{i}(x)], \quad y^{i}(x) \coloneq F(\xi^{i},x), \quad \phi = [\xi^{i},x]
	\end{equation}

	\begin{equation}\label{eq:def-y-diag}
		y^{i}(x)_{i} = -\sum_{j \in \nu} \frac{\kappa_{ij}}{Z_{i}}(x_{j} - x_{i}) \coloneq (\Delta_{\eta,\kappa} x)_{i} \in \R^{d}.
	\end{equation}

	\begin{equation}\label{eq:ass-g}
		\xi^{ij} \coloneq (\xi^{i})^{-1} \xi^{j} \in \Xi
	\end{equation}

	\begin{proposition}[consistent locally Laplacian operators are gauged Laplacians]\label{prop:loc-lpl}
		Let $\mc{F} : \Gamma \to \Gamma$ be a locally Laplacian gauge invariant data transformation, such that the associated gauges $\xi^{ij}$ satisfy the relations
		\begin{equation}\label{eq:intergrability}
			\xi^{ij}_{i} = \xi^{ij}_{j}, \qquad \forall i,j \in \nu.
		\end{equation}
		Then there is a gauge invariant graph voltage $\omega \in \Omega$ such that
		\begin{equation}\label{eq:mcF=Lchi}
			\mc{F} = L_{\chi}, \qquad \chi = (\eta,\kappa,\omega),
		\end{equation}
		and for all $i,j,k \in \nu, j > k$, we have
		\begin{equation}\label{eq:def-alphai}
			\omega = [\xi^{i},\alpha^{i}], \qquad
			\alpha^{i}_{jk} = \xi^{ij}_{j} \xi^{ji}_{k} \in \SO(d).
		\end{equation}
	\end{proposition}

	\begin{proof}
		We show that $\omega$ defines a gauge invariant voltage, which amounts to verifying that the following consistency conditions hold
		\begin{equation}
			\alpha_{jk}^{i} = (\alpha_{kj}^{i})^{-1}, \quad \text{and} \quad
			\alpha_{jk}^{m} = \xi^{mi}_{j} \alpha^{i}_{jk} \xi^{im}_{k}
		\end{equation}
		for all $i,j,k,m \in \nu$. The second equation one holds since
		\begin{equation}
			\xi^{mi}_{j} \alpha^{i}_{jk} \xi^{im}_{k}
			=
			\underbrace{\xi^{mi}_{j} \xi^{ij}_{j}}_{\overset{\eqref{eq:ass-g}}{=}\xi^{mj}_{j}}
			\underbrace{\xi^{ji}_{k} \xi^{im}_{k}}_{\overset{\eqref{eq:ass-g}}{=}\xi^{jm}_{k}}
			= \xi^{mj}_{j} \xi^{jm}_{k}
			\overset{\eqref{eq:def-alphai}}{=} \alpha^{m}_{jk}.
		\end{equation}
		For the first condition, we compute
		\begin{equation}
			\alpha^{i}_{jk}\alpha^{i}_{kj}
			=
			\xi^{ij}_{j}
			\underbrace{
				\xi^{ji}_{k}\xi^{ik}_{k}
			}_{\overset{\eqref{eq:ass-g}}{=} \xi^{jk}_{k}}
			\xi^{ki}_{j}
			=
			\xi^{ij}_{j}
			\underbrace{
				\xi^{jk}_{k}
			}_{\overset{\eqref{eq:intergrability}}{=} \xi^{jk}_{j}}
			\xi^{ki}_{j}
			=
			\xi^{ij}_{j}\xi^{jk}_{j}\xi^{ki}_{j}
			\overset{\eqref{eq:ass-g}}{=}
			\II.
		\end{equation}
		This shows that the definition \eqref{eq:def-alphai} is consistent and defines a graph voltage.

		We proceed in order to prove \eqref{eq:mcF=Lchi}.
		Since both transformations are gauge invariant, we can check equality in any gauge, which we choose here to be $\xi^{i}$ for a fixed choice of $i \in \nu$. Evaluating both operators on $\phi = [\xi^{i},x]$ yields
		\begin{equation}
			\mc{F}([\xi^{i},x]) = [\xi^{i},y^{i}(x)], \qquad
			L_{\chi}([\xi^{i},x]) = [\xi^{i},\Delta_{\eta,\kappa,\alpha^{i}}x],
		\end{equation}
	 	where $y^{i}(x)$ is defined by \eqref{eq:def-yi}.
		We thus have to verify that
		\begin{equation}\label{eq:proof-show-Delta-y-i}\Delta_{\eta,\kappa,\alpha^{i}}x = y^{i}(x).
		\end{equation}
		We separately check equality at all nodes $j \in \nu$ and first calculate $y^{i}(x)_j$. Due to gauge invariance, the equation
		\begin{equation}
			y^{i}(x) = F(\xi^{i},x)
			\overset{
				\substack{
					\eqref{eq:consis} \\ \eqref{eq:ass-g}
				}
			}{=}
			\xi^{ij} F(\xi^{j},\xi^{ji}x)
			\overset{
				\eqref{eq:def-yi}
			}{=}
			\xi^{ij} y^{j}(\xi^{ji}x).
		\end{equation}
		holds. In particular, we have
		\begin{equation}
			y^{i}(x)_{j} = \xi^{ij}_{j} y^{j}(\xi^{ji}x)_{j}, \quad \forall j \in \nu.
		\end{equation}
		On the other hand, relation \eqref{eq:def-y-diag} yields
		\begin{equation}
			y^{j}(\xi^{ji} x)_{j} = (\Delta_{\eta,\kappa} \xi^{ji}x)_{j}.
		\end{equation}
		By combining both equations, we obtain
		\begin{equation}\label{eq:proof-F-L-aux-1}
			y^{i}(x)_{j}
			= \xi^{ij}_{j}y^{j}(\xi^{ji}x)_j
			= \xi^{ij}_{j}(\Delta_{\eta,\kappa} \xi^{ji}x )_{j}
		\end{equation}
		which expands to
		\begin{equation}
			\xi^{ij}_{j}(\Delta_{\eta,\kappa} \xi^{ji}x )_{j}
			\overset{\eqref{eq:def-y-diag}}{=}
			- \xi^{ij}_{j} \sum_{k\in \nu} \frac{\kappa_{jk}}{Z_{j}}(\xi^{ji}_{k} x_{k} - \xi^{ji}_{j} x_{j})
			=
			- \sum_{k \in \nu} \frac{\kappa_{jk}}{Z_{j}}(\xi^{ij}_{j}\xi^{ji}_{k}x_{k} - x_{j})
		\end{equation}
		using $\xi^{ij}_{j}\xi^{ji}_{j} = \II$ in the last step. Applying the definition \eqref{eq:def-alphai} to the last expression gives
		\begin{equation}\label{eq:proof-F-L-aux-2}
			\xi^{ij}_{j}(\Delta_{\eta,\kappa} \xi^{ji}x)_{j}
			= -\sum_{k \in \nu} \frac{\kappa_{jk}}{Z_{j}}(\alpha_{jk}^{i} x_{k} - x_j)
			\overset{\eqref{eq:def-lalpha}}{=} (\Delta_{\eta,\kappa,\alpha^{i}}x)_{j}.
		\end{equation}
		Hence \eqref{eq:proof-F-L-aux-1}--\eqref{eq:proof-F-L-aux-2} show that
		\begin{equation}
			y^{i}(x)_{j} = (\Delta_{\eta,\kappa,\alpha^{i}}x)_{j}, \quad \forall j \in \nu.
		\end{equation}
		We conclude that Equation \eqref{eq:proof-show-Delta-y-i} holds and $\mc{F} = L_{\chi}$ with $\chi=(\eta,\kappa,\omega)$.

		We finish the proof by demonstrating that the equation $\mc{F} = L_{\chi}$ is compatible with $\mc{F}$ being locally Laplacian, i.e. the assumptions of the Proposition are not violated. By virtue of the definitions \eqref{eq:ass-g} and \eqref{eq:def-alphai}, we find $\alpha^{i}_{ij} =\II$ for all edges $i>j$, such that
			\begin{equation}
				y^{i}(x)_{i}
				= - \sum_{j\in\nu}\frac{\kappa_{ij}}{Z_{i}}(\alpha^{i}_{ij} x_{j} - x_{i})
				= - \sum_{j \in \nu} \frac{\kappa_{ij}}{Z_{j}}(x_{j} - x_{i})
				= (\Delta_{\eta,\kappa} x)_{i}, \quad \forall i \in \nu,
			\end{equation}
		which is Equation \eqref{eq:def-y-diag}.
	\end{proof}

\end{document}
